\documentclass[UTF8,12pt]{ctexart}
\usepackage[a4paper,margin=24mm]{geometry}
\usepackage{amsmath,amssymb,bm,graphicx,adjustbox}
\newcommand{\fitmath}[1]{\begin{adjustbox}{max width=\linewidth}$\displaystyle #1$\end{adjustbox}}
\setlength{\parindent}{0pt}
\setlength{\parskip}{.7em}
\sloppy
\allowdisplaybreaks
\begin{document}
\section*{2014 年考研数学三 · 参考答案}
\subsection*{原 PDF 第 30 页}

\subsection*{2014 年真题参考答案}

\subsection*{一、选择题}

(1) A。 (2) C。 (3) D。 (4) D。 (5) B。 (6) A。 (7) B。 (8) C。


\subsection*{二、填空题}

(9) \(\displaystyle 20-Q\)。 (10) \(\displaystyle \dfrac{\displaystyle 3}{\displaystyle 2}-\ln2\)。 (11) \(\displaystyle \dfrac{\displaystyle 1}{\displaystyle 2}\)。 (12) \(\displaystyle \dfrac{\displaystyle e-1}{\displaystyle 2}\)。 (13) \(\displaystyle [-2,\allowbreak 2]\)。 (14) \(\displaystyle \dfrac{\displaystyle 2}{\displaystyle 5n}\)。


\subsection*{三、解答题}

(15) \(\displaystyle \dfrac{\displaystyle 1}{\displaystyle 2}\)。


(16) \(\displaystyle -\dfrac{\displaystyle 3}{\displaystyle 4}\)。


(17) \(\displaystyle f(u)=\dfrac{\displaystyle 1}{\displaystyle 16}e^{4u}-\dfrac{\displaystyle 1}{\displaystyle 4}u-\dfrac{\displaystyle 1}{\displaystyle 16}\)。


(18) 和函数 \(\displaystyle s(x)=\dfrac{\displaystyle 3-x}{\displaystyle (1-x)^3}\)，收敛域为 \(\displaystyle (-1,\allowbreak 1)\)。


(19) 证明略。


(20)（I）\(\displaystyle (-1,\allowbreak 2,\allowbreak 3,\allowbreak 1)^{\mathrm T}\)。（II）\(\displaystyle B=\begin{pmatrix}2-k_1&6-k_2&-1-k_3\\-1+2k_1&-3+2k_2&1+2k_3\\-1+3k_1&-4+3k_2&1+3k_3\\k_1&k_2&k_3\end{pmatrix}\)，其中 \(\displaystyle k_1,\allowbreak k_2,\allowbreak k_3\) 为任意常数。


(21) 证明略。


(22)（I）\(\displaystyle F_Y(y)=\begin{cases}0,\allowbreak &y<0,\allowbreak \\\dfrac{\displaystyle 3}{\displaystyle 4}y,\allowbreak &0\leq y<1,\allowbreak \\\dfrac{\displaystyle 1}{\displaystyle 2}+\dfrac{\displaystyle 1}{\displaystyle 4}y,\allowbreak &1\leq y<2,\allowbreak \\1,\allowbreak &y\geq2.\end{cases}\)（II）\(\displaystyle \dfrac{\displaystyle 3}{\displaystyle 4}\)。


(23)（I）\(\displaystyle (X,\allowbreak Y)\) 的概率分布为 \(\displaystyle \begin{array}{c|cc}Y\backslash X&0&1\\\hline 0&\dfrac{\displaystyle 2}{\displaystyle 9}&\dfrac{\displaystyle 1}{\displaystyle 9}\\1&\dfrac{\displaystyle 1}{\displaystyle 9}&\dfrac{\displaystyle 5}{\displaystyle 9}\end{array}\)。（II）\(\displaystyle \dfrac{\displaystyle 4}{\displaystyle 9}\)。


\clearpage
\end{document}
